% ============================================================
%  teacher_guide.tex — Robot World: Meet the Machines That Help Us
%  Parent & Teacher Guide
% ============================================================

\namedchapter{Parent \& Teacher Guide}
\addcontentsline{toc}{chapter}{Parent \& Teacher Guide}

\begin{tcolorbox}[enhanced,colback=SunYellow!12,colframe=RoboGreen!60!black,arc=6pt,boxrule=1pt,left=10pt,right=10pt,top=8pt,bottom=8pt]
\textbf{Cognitive load note:} This book is more demanding than Books 1 and 2. For learners aged 8--9, pre-teach vocabulary from ROS, machine learning, AI ethics, swarm robotics, and autonomy.\\
\textbf{AI ethics and companion robots prompts:} Use guided reflections such as ``What do you think it means for a robot to be a friend? How is that different from a human friend?''\\
\textbf{Illustration density note:} Chapters 4 and 10 are relatively illustration-light. Supplement with printed infographics or whiteboard drawings during instruction.
\end{tcolorbox}

This guide offers chapter-by-chapter discussion questions, extension
activities, and curriculum links to help adults support young readers.
Use it alongside classroom lessons, home reading, or after-school STEM clubs.

\sectionrule

% --- How to use this book ---
\section*{How to Use This Book}

\begin{itemize}
  \item \textbf{Read-aloud or independent:} Chapters work well read aloud (ages 7--9)
    or independently (Ages 10--13).
  \item (Less suitable for independent readers under 10; adult guidance recommended for younger readers)
  \item \textbf{One chapter per session:} Each chapter is self-contained and takes
    approximately 15--20 minutes to read.
  \item \textbf{Pair with hands-on:} Every chapter includes a Try It or Experiment
    activity. Let children try before moving on.
  \item \textbf{Use the Quick Check:} The three-question review at the end of each
    chapter is ideal for checking comprehension before discussion.
  \item \textbf{Revisit the Glossary:} Encourage children to look up bold terms
    in the back-matter glossary.
\end{itemize}

\sectionrule

% --- Chapter-by-chapter ---
\section*{Chapter Notes}

\subsection*{Chapter 1: What Is a Robot?}
\begin{itemize}
  \item \textbf{Key concept:} Sense\,--\,think\,--\,act loop.
  \item \textbf{Discussion:}
    \begin{enumerate}
      \item Ask children to classify household items as robots or simple machines.
      \item What makes a robot vacuum different from a regular vacuum?
      \item Can something be a machine but not a robot? Give examples.
    \end{enumerate}
  \item \textbf{Activity:} Create a ``robot or not?'' poster collage from magazines or printed images.
  \item \textbf{Extension:} Research Karel \v{C}apek's play \emph{R.U.R.} and draw a timeline of robots in fiction.
  \item \textbf{Curriculum link:} Science --- forces and mechanisms; Design \& Technology --- understanding systems.
\end{itemize}

\subsection*{Chapter 2: Robot Bodies}
\begin{itemize}
  \item \textbf{Key concept:} Motors, joints, grippers, materials, and energy.
  \item \textbf{Discussion:}
    \begin{enumerate}
      \item Why do different robots need different body shapes?
      \item What would happen if a drone used wheels instead of propellers?
      \item How is a robot arm's joint like your elbow?
    \end{enumerate}
  \item \textbf{Activity:} Build a cardboard gripper (lolly-stick experiment in book).
  \item \textbf{Extension:} Compare the materials in a bicycle, a drone, and a submarine. Why are they different?
  \item \textbf{Curriculum link:} D\&T --- mechanisms and structures; Science --- properties of materials.
\end{itemize}

\subsection*{Chapter 3: How Robots Sense the World}
\begin{itemize}
  \item \textbf{Key concept:} Sensors and sensor fusion.
  \item \textbf{Discussion:}
    \begin{enumerate}
      \item Which human senses are hardest to replicate electronically?
      \item Why might a self-driving car need both cameras and lidar?
      \item What happens if one sensor breaks --- can the robot still work?
    \end{enumerate}
  \item \textbf{Activity:} Blindfold obstacle course to experience limited sensing.
  \item \textbf{Extension:} Research how bats use echolocation and compare to ultrasonic sensors.
  \item \textbf{Curriculum link:} Science --- light and sound; Computing --- inputs and outputs.
\end{itemize}

\subsection*{Chapter 4: The Robot Brain}
\begin{itemize}
  \item \textbf{Key concept:} Programmes, algorithms, loops, conditions, and machine learning.
  \item \textbf{Discussion:}
    \begin{enumerate}
      \item Why must robot instructions be very precise?
      \item What everyday activities can you describe as an algorithm?
      \item How is machine learning different from a fixed programme?
    \end{enumerate}
  \item \textbf{Activity:} Programme a partner: give step-by-step voice commands to navigate a room.
  \item \textbf{Extension:} Try Scratch or Open Roberta to write a simple robot programme on screen.
  \item \textbf{Curriculum link:} Computing --- algorithms and programming; Maths --- sequencing.
\end{itemize}

\subsection*{Chapter 5: Factory Robots}
\begin{itemize}
  \item \textbf{Key concept:} Assembly lines, repetition, cobots, and warehouse automation.
  \item \textbf{Discussion:}
    \begin{enumerate}
      \item What jobs are best for robots vs.\ humans in a factory?
      \item Would you rather work with a cobot or do the task alone? Why?
      \item How do warehouse robots change the way we shop online?
    \end{enumerate}
  \item \textbf{Activity:} Mini assembly-line relay race (book's Try It).
  \item \textbf{Extension:} Watch a short video of car factory robots. Count how many different tasks they do.
  \item \textbf{Curriculum link:} History --- Industrial Revolution; PSHE --- work and careers; Maths --- time and efficiency.
\end{itemize}

\subsection*{Chapter 6: Robots at Home}
\begin{itemize}
  \item \textbf{Key concept:} Autonomy, mapping, charging cycles, and companion robots.
  \item \textbf{Discussion:}
    \begin{enumerate}
      \item What chores would you give a home robot?
      \item Should voice assistants always be listening? Why or why not?
      \item Could a robot vacuum get ``confused''? How?
    \end{enumerate}
  \item \textbf{Activity:} Map a room blindfolded (book's experiment).
  \item \textbf{Extension:} Draw a floor plan of your home and mark the best path for a robot vacuum.
  \item \textbf{Curriculum link:} Computing --- data and information; Geography --- mapping and plans.
\end{itemize}

\subsection*{Chapter 7: Robots in Space and Deep Sea}
\begin{itemize}
  \item \textbf{Key concept:} Extreme environments, signal delay, autonomy.
  \item \textbf{Discussion:}
    \begin{enumerate}
      \item Why is signal delay a problem for controlling Mars rovers?
      \item What would happen if an undersea cable broke --- how would engineers fix it?
      \item Why are rescue robots better than humans in some disaster situations?
    \end{enumerate}
  \item \textbf{Activity:} Walkie-talkie delay game (book's experiment).
  \item \textbf{Extension:} Research Perseverance's instruments and create a labelled poster.
  \item \textbf{Curriculum link:} Science --- Earth and space; Geography --- oceans and mapping.
\end{itemize}

\subsection*{Chapter 8: Medical Robots}
\begin{itemize}
  \item \textbf{Key concept:} Precision, telesurgery, exoskeletons, and micro-bots.
  \item \textbf{Discussion:}
    \begin{enumerate}
      \item Would you trust a robot to perform surgery? Why or why not?
      \item How might exoskeletons help people with disabilities?
      \item What are the risks of nano-bots inside the body?
    \end{enumerate}
  \item \textbf{Activity:} Steady-hand game (book's Try It).
  \item \textbf{Extension:} Compare the size of the da Vinci robot's instruments to everyday objects.
  \item \textbf{Curriculum link:} Science --- the human body; PSHE --- health and wellbeing; D\&T --- precision.
\end{itemize}

\subsection*{Chapter 9: Robot Friends}
\begin{itemize}
  \item \textbf{Key concept:} Social robots, emotion detection, ethics, and privacy.
  \item \textbf{Discussion:}
    \begin{enumerate}
      \item Can robots really be friends? Explore feelings vs.\ programming.
      \item Should we treat robots kindly even though they cannot feel?
      \item If a companion robot records your voice, who owns that data?
    \end{enumerate}
  \item \textbf{Activity:} Design a companion robot (book's experiment).
  \item \textbf{Extension:} Hold a class debate: ``Companion robots do more harm than good.''
  \item \textbf{Curriculum link:} PSHE --- relationships, empathy, and digital citizenship; Computing --- AI.
\end{itemize}

\subsection*{Chapter 10: The Future of Robots}
\begin{itemize}
  \item \textbf{Key concept:} Soft robots, swarms, machine learning, bio-inspiration, safety, and ethics.
  \item \textbf{Discussion:}
    \begin{enumerate}
      \item What problems would you like robots to solve in the future?
      \item How can we make sure robots are fair and safe for everyone?
      \item What jobs should robots \emph{never} do? Why?
    \end{enumerate}
  \item \textbf{Activity:} Invent-a-robot challenge (book's Try It).
  \item \textbf{Extension:} Write a short story set 50 years in the future featuring a robot helper.
  \item \textbf{Curriculum link:} D\&T --- design process; Citizenship --- technology and society; English --- creative writing.
\end{itemize}

\sectionrule

% --- Cross-curricular projects ---
\section*{Cross-Curricular Project Ideas}

\begin{enumerate}
  \item \textbf{Robot Museum (Art + Science):} Each child designs and builds a robot
    model from recycled materials, writes a museum-style label, and presents it to the class.
  \item \textbf{Robot Timeline (History + Computing):} Create a classroom wall timeline
    from the first automata (ancient Greece) to modern-day AI robots.
  \item \textbf{Robot Safety Code (PSHE + English):} Write a ``Robot Rules'' charter
    covering privacy, fairness, and safety. Display it in the classroom.
  \item \textbf{Programmable Robot Day (Computing + Maths):} Use Bee-Bots, mBots,
    or Scratch to programme a robot to navigate a maze.
  \item \textbf{Debate Day (English + Citizenship):} ``Should robots replace human
    workers?'' --- formal debate with research, evidence, and rebuttals.
\end{enumerate}

\sectionrule

% --- General tips ---
\section*{General Teaching Tips}

\begin{itemize}
  \item Encourage children to look for robots in their daily lives and keep a ``robot diary.''
  \item Pair chapters with free online robot simulators (e.g.\ Open Roberta, Scratch with mBot, CoderZ).
  \item Use the glossary as a spelling/vocabulary challenge --- award points for using terms in sentences.
  \item The quiz can be done individually, as a team quiz-show format, or as a take-home challenge.
  \item Connect back to Book 1 (Internet) and Book 2 (Smartphones) whenever relevant ---
    robots use networks, processors, and sensors that children already understand.
  \item Preview Book 4 (\emph{When Machines Start Thinking}) by discussing the machine learning
    section in Chapter 4.
\end{itemize}

\sectionrule

% --- Further reading ---
\section*{Further Reading \& Resources}

\begin{itemize}
  \item \textbf{Books:}
    \begin{itemize}
      \item \emph{Robots: Meet the Machines of the Future} (DK)
      \item \emph{How to Build a Robot} (Lonely Planet Kids)
      \item \emph{The Way Things Work Now} by David Macaulay
    \end{itemize}
  \item \textbf{Websites:}
    \begin{itemize}
      \item NASA Mars Exploration --- \texttt{mars.nasa.gov}
      \item Scratch Programming --- \texttt{scratch.mit.edu}
      \item Open Roberta Lab --- \texttt{lab.open-roberta.org}
    \end{itemize}
  \item \textbf{Videos:}
    \begin{itemize}
      \item Boston Dynamics YouTube channel (robot demos)
      \item How It's Made (factory robot episodes)
      \item BBC Bitesize: Computing --- Robots and AI
    \end{itemize}
\end{itemize}

\sectionrule

% --- Series info ---
\begin{tcolorbox}[enhanced,colback=LightGray,colframe=RoboGreen!40,arc=6pt,boxrule=0.8pt,
  left=10pt,right=10pt,top=8pt,bottom=8pt]
\textbf{Series:} The Secret Life of Technology\\
\textbf{Book 1:} The Secret Life of the Internet: A Journey Through Wires, Waves, and Wonders \checkmark\\
\textbf{Book 2:} Inside a Smartphone: The Tiny World in Your Hand \checkmark\\
\textbf{Book 3:} Robot World: Meet the Machines That Help Us \checkmark\\
\textbf{Coming next:} When Machines Start Thinking: A Kid's Guide to AI
\end{tcolorbox}
